
\documentclass{article}
\usepackage{colortbl}
\usepackage{makecell}
\usepackage{multirow}
\usepackage{supertabular}

\begin{document}

\newcounter{utterance}

\centering \large Interaction Transcript for game `referencegame', experiment `random\_grids', episode 14 with reflexion{-}llama{-}3.1{-}8b{-}instruct.
\vspace{24pt}

{ \footnotesize  \setcounter{utterance}{1}
\setlength{\tabcolsep}{0pt}
\begin{supertabular}{c@{$\;$}|p{.15\linewidth}@{}p{.15\linewidth}p{.15\linewidth}p{.15\linewidth}p{.15\linewidth}p{.15\linewidth}}
    \# & $\;$A & \multicolumn{4}{c}{Game Master} & $\;\:$B\\
    \hline

    \theutterance \stepcounter{utterance}  
    & & \multicolumn{4}{p{0.6\linewidth}}{
        \cellcolor[rgb]{0.9,0.9,0.9}{
            \makecell[{{p{\linewidth}}}]{
                \texttt{\tiny{[P1$\langle$GM]}}
                \texttt{You have attempted this same task before and did not succeed. These are the plans you wrote after those attempts.} \\
\texttt{Trial \#1: New plan: I failed to distinguish the target grid from the distractor grids because my description was too general. The target grid has a row with three X's in the first, third, and fifth columns, but this description also applies to the first row of the distractor grid 2. Next time, I will look for a more specific feature that is unique to the target grid. I will examine the columns and rows more carefully to find a distinguishing characteristic, such as a specific pattern of X's and empty cells. For example, I might look for a row or column that has a unique combination of X's and empty cells, or a specific arrangement of X's that is not found in the distractor grids.} \\
\\ 
\texttt{Trial \#2: New plan: I failed to distinguish the target grid from the distractor grids because my description was too general. The target grid has a row with three X's in the first, third, and fifth columns, but this description also applies to the first row of the distractor grid 2. Next time, I will look for a more specific feature that is unique to the target grid. I will examine the columns and rows more carefully to find a distinguishing characteristic, such as a specific pattern of X's and empty cells. For example, I might look for a row or column that has a unique combination of X's and empty cells, or a specific arrangement of X's that is not found in the distractor grids. } \\
\\ 
\texttt{In particular, I will pay closer attention to the middle row of the target grid, which has a unique combination of X's and empty cells. I will try to describe this row in a way that distinguishes it from the corresponding rows in the distractor grids. For example, I might say that the target grid has a row with an X in the second column and an X in the fourth column, but not in the first or third columns.} \\
\\ 
\\ 
\texttt{You are given three grids, where each of them is 5 by 5 in size.} \\
\texttt{Grids have empty cells marked with "▢" and filled cells marked with "X".} \\
\texttt{Your task is to generate a referring expression that best describes the target grid while distinguishing it from the two other distractor grids.} \\
\texttt{The first grid is the target grid, and the following two grids are the distractors.} \\
\\ 
\texttt{Target grid:} \\
\\ 
\texttt{□ X X □ X} \\
\texttt{X □ X X □} \\
\texttt{□ X □ X X} \\
\texttt{X □ X □ X} \\
\texttt{□ X X □ □} \\
\\ 
\texttt{Distractor grid 1:} \\
\\ 
\texttt{□ X □ X □} \\
\texttt{X □ X □ X} \\
\texttt{□ X □ X □} \\
\texttt{X □ X □ X} \\
\texttt{□ X □ X □} \\
\\ 
\texttt{Distractor grid 2:} \\
\\ 
\texttt{□ □ X X □} \\
\texttt{X □ □ X X} \\
\texttt{□ X X □ X} \\
\texttt{X X □ □ X} \\
\texttt{□ X X □ □} \\
\\ 
\texttt{Instruction: Describe the target grid.} \\
\texttt{Generate the referring expression starting with the tag "Expression: " for the given target grid. Omit any other text.} \\
            }
        }
    }
    & \\ \\

    \theutterance \stepcounter{utterance}  
    & \multicolumn{4}{p{0.6\linewidth}}{
        \cellcolor[rgb]{0.9,0.9,0.9}{
            \makecell[{{p{\linewidth}}}]{
                \texttt{\tiny{[P1$\rangle$GM]}}
                \texttt{Expression: The grid that has a row with an X in the first column, an X in the third column, and an X in the fifth column, and a row with an X in the second column and an X in the fourth column, but not in the first or third columns.} \\
            }
        }
    }
    & & \\ \\

    \theutterance \stepcounter{utterance}  
    & & & \multicolumn{2}{p{0.3\linewidth}}{
        \cellcolor[rgb]{0.9,0.9,0.9}{
            \makecell[{{p{\linewidth}}}]{
                \texttt{\tiny{[GM$|$GM]}}
                \texttt{Expression: The grid that has a row with an X in the first column, an X in the third column, and an X in the fifth column, and a row with an X in the second column and an X in the fourth column, but not in the first or third columns.} \\
            }
        }
    }
    & & \\ \\

    \theutterance \stepcounter{utterance}  
    & & \multicolumn{4}{p{0.6\linewidth}}{
        \cellcolor[rgb]{0.9,0.9,0.9}{
            \makecell[{{p{\linewidth}}}]{
                \texttt{\tiny{[GM$\rangle$P2]}}
                \texttt{You have attempted this same task before and did not succeed. These are the plans you wrote after those attempts.} \\
\texttt{Trial \#1: New plan: I failed to identify the correct grid because I didn't carefully examine the expression. The expression states that the grid has a row with three X's in the first, third, and fifth columns. However, I didn't check if the expression is referring to a column or a row. Next time, I will carefully read the expression and check if it's referring to a row or a column. I will also examine each grid more closely to see if it matches the expression.} \\
\\ 
\texttt{Trial \#2: New plan: I failed to identify the correct grid because I didn't carefully examine the expression. The expression states that the grid has a row with three X's in the first, third, and fifth columns and a row with three X's in the first, second, and fifth columns. However, I didn't check if the expression is referring to a column or a row. Next time, I will carefully read the expression and check if it's referring to a row or a column. I will also examine each grid more closely to see if it matches the expression. Specifically, I will look for rows that match the given pattern, and I will also check if the pattern is referring to columns. If I find a row that matches the pattern, I will check if the other row also matches the pattern. If I find a grid that matches both patterns, I will select that grid as the answer.} \\
\\ 
\\ 
\texttt{You are given three grids, where each of them is 5 by 5 in size.} \\
\texttt{Grids have empty cells marked with "▢" and filled cells marked with "X".} \\
\texttt{You are also given a referring expression that describes one of the given grids.} \\
\texttt{Your task is to select the grid that matches the given referring expression.} \\
\texttt{Generate only the number (in text) of the grid that matches the given expression by selecting first, second, or third.} \\
\\ 
\texttt{First:} \\
\\ 
\texttt{□ □ X X □} \\
\texttt{X □ □ X X} \\
\texttt{□ X X □ X} \\
\texttt{X X □ □ X} \\
\texttt{□ X X □ □} \\
\\ 
\texttt{Second:} \\
\\ 
\texttt{□ X □ X □} \\
\texttt{X □ X □ X} \\
\texttt{□ X □ X □} \\
\texttt{X □ X □ X} \\
\texttt{□ X □ X □} \\
\\ 
\texttt{Third:} \\
\\ 
\texttt{□ X X □ X} \\
\texttt{X □ X X □} \\
\texttt{□ X □ X X} \\
\texttt{X □ X □ X} \\
\texttt{□ X X □ □} \\
\\ 
\texttt{Expression: The grid that has a row with an X in the first column, an X in the third column, and an X in the fifth column, and a row with an X in the second column and an X in the fourth column, but not in the first or third columns.} \\
\texttt{Question: Which grid does the expression refer to?} \\
\texttt{Start with the tag "Answer: ", followed by your selection. Omit any other text.} \\
            }
        }
    }
    & \\ \\

    \theutterance \stepcounter{utterance}  
    & & & \multicolumn{4}{p{0.6\linewidth}}{
        \cellcolor[rgb]{0.9,0.9,0.9}{
            \makecell[{{p{\linewidth}}}]{
                \texttt{\tiny{[GM$\langle$P2]}}
                \texttt{Answer: third} \\
            }
        }
    }
     \\ \\

    \theutterance \stepcounter{utterance}  
    & & & \multicolumn{2}{p{0.3\linewidth}}{
        \cellcolor[rgb]{0.9,0.9,0.9}{
            \makecell[{{p{\linewidth}}}]{
                \texttt{\tiny{[GM$|$GM]}}
                \texttt{third} \\
            }
        }
    }
    & & \\ \\

\end{supertabular}
}

\end{document}
